\documentclass[11pt,a4paper]{article}
\usepackage[a4paper,margin=25mm,headheight=15pt]{geometry}
\usepackage{fontspec}
\defaultfontfeatures{Ligatures=TeX}
% Load TeX Live fonts by filename: arXiv has limited fontconfig registration.
\setmainfont{lmroman10-regular.otf}[
BoldFont=lmroman10-bold.otf,ItalicFont=lmroman10-italic.otf,
BoldItalicFont=lmroman10-bolditalic.otf]
\setmonofont[Scale=0.88]{lmmono10-regular.otf}
\usepackage{amsmath,unicode-math,microtype,enumitem,titlesec,fancyhdr,xurl}
\setmathfont{latinmodern-math.otf}
\usepackage[hidelinks,unicode]{hyperref}
\usepackage{hyperxmp}
\hypersetup{pdftitle={ROS3-Limassol: An OpenPELAOS Prototype for Telemetry Liveness Monitoring},
pdfauthor={Charalampos Nadiotis},
pdfsubject={Preliminary technical report; civilian-use licence; military use prohibited; initial idea credited to Christos Paparistodemou},
pdfkeywords={robotics, telemetry, liveness monitoring, ROS 2, intermittent connectivity, OpenPELAOS, civilian use, no military use},
pdflang={en},
pdfcopyright={Copyright 2026 Charalampos Nadiotis. Custom Civilian-Use Licence, Version 1.0: military use prohibited. Full terms in this manuscript. Initial project idea: Christos Paparistodemou.}}
\linespread{1.02}
\setlength{\parindent}{0pt}
\setlength{\parskip}{4pt}
\setlength{\emergencystretch}{2em}
\setlist{nosep,leftmargin=1.4em}
\titleformat{\section}{\large\bfseries}{\thesection}{0.65em}{}
\titleformat{\subsection}{\normalsize\bfseries}{\thesubsection}{0.65em}{}
\titlespacing*{\section}{0pt}{12pt}{5pt}
\titlespacing*{\subsection}{0pt}{8pt}{3pt}
\pagestyle{fancy}
\fancyhf{}
\fancyhead[L]{\small OpenPELAOS / ROS3-Limassol}
\fancyhead[R]{\small Preliminary technical report}
\fancyfoot[L]{\scriptsize\textcopyright\ 2026 C. Nadiotis\quad\hyperref[sec:licence]{Civilian-Use Licence - No military use}}
\fancyfoot[R]{\small\thepage}
\fancypagestyle{plain}{\fancyhf{}\renewcommand{\headrulewidth}{0pt}
\fancyfoot[L]{\scriptsize\textcopyright\ 2026 C. Nadiotis\quad\hyperref[sec:licence]{Civilian-Use Licence - No military use}}
\fancyfoot[R]{\small\thepage}}
\renewcommand{\headrulewidth}{0.3pt}
\newcommand{\code}[1]{\texttt{#1}}
\newcommand{\healthy}{\mathrm{HEALTHY}}
\newcommand{\stale}{\mathrm{STALE}}
\newcommand{\waiting}{\mathrm{WAITING}}
\begin{document}
\thispagestyle{plain}
\begin{center}
{\LARGE\bfseries ROS3-Limassol: An OpenPELAOS Prototype\\[3pt]
for Telemetry Liveness Monitoring}\\[8pt]
{\large Intermittently Connected Robotic Systems}\\[11pt]
Charalampos Nadiotis\\[3pt]
OpenPELAOS, independent open-source project\\[5pt]
{\small Initial project idea: Christos Paparistodemou}\\[5pt]
{\small ResearchGate preprint edition: 25 October 2026}\\[3pt]
{\small Original experiment: 23 October 2026, software v0.1.0}\\[5pt]
{\footnotesize\textcopyright\ 2026 Charalampos Nadiotis.\quad
\hyperref[sec:licence]{Custom Civilian-Use Licence, v1.0.}}\\[2pt]
{\footnotesize\textbf{Military use is prohibited. Full licence terms appear in this report.}}
\end{center}

\begin{abstract}
Robotic telemetry can stop arriving because of communication disruption, publisher
failure, or upstream data interruption. A receiver therefore needs an explicit
rule for interpreting the absence of valid messages. This report describes
ROS3-Limassol, an OpenPELAOS development prototype comprising a transport-independent
Python core, a synthetic unmanned surface vehicle telemetry source, and an optional
ROS~2 adapter. The observer evaluates elapsed receiver time since the last accepted
message and distinguishes waiting, healthy, and stale states. Schema validation,
session identifiers, and increasing sequence numbers prevent several malformed
or replayed messages from refreshing that timer. A deterministic experiment
with a three-second blackout reproduces stale detection and recovery, accepting
51 of 81 generated samples. In the original version-0.1.0 validation, eighteen core tests pass; one ROS graph test is
explicitly skipped because ROS~2 is unavailable in the authoring environment.
The evidence establishes software behavior under a specified synthetic scenario.
Real transport timing, field performance, and research novelty remain unestablished.
\end{abstract}
{\small\textbf{Keywords:} ROS~2; field robotics; telemetry; liveness monitoring;
intermittent connectivity; reproducibility.}

\section{Introduction}
A robotics operator may continue to see an apparently valid position or battery
reading after the source has stopped delivering usable updates. The most recent
value and the current availability of its source are different observations.
ROS3-Limassol begins with the task of making arrival liveness explicit, before
attempting navigation, mission execution, or broader autonomy.

The project is developed under OpenPELAOS, with the initial idea credited to
Christos Paparistodemou. ROS3-Limassol is the project identifier; version~0.1.0
implements a toolkit that can use ROS~2. The immediate outcome is a small
implementation that can be inspected and exercised without a robot or a ROS
installation. Its intended ROS mode adds interprocess communication around the
same core logic. This is an independent project, rather than an official
ROS distribution.

The motivating question is whether a bounded receiver-side observer can report
the absence and resumption of accepted telemetry without depending on synchronized
sender and receiver clocks. The present contribution is an implementation and
deterministic verification baseline \cite{artifact}. It does not introduce a new
failure-detection algorithm or establish an advantage over existing robotics
software. Its value at this stage is a precise contract, reproducibility, and
an explicit boundary between verified behavior and planned evaluation.

\section{Background and technical basis}
ROS~2 provides an established modular robotics software architecture.
Macenski et al.\ describe its design and deployment context
\cite{macenski2022}. Reusing that infrastructure allows the prototype to focus on
application-level observations. The adapter follows the Python publisher/subscriber
and \code{ament\_python} conventions in the official documentation
\cite{ros_python}.

ROS~2 quality-of-service settings distinguish history, reliability, and durability,
and communicating endpoints require compatible policies \cite{ros_qos}.
These policies do not, by themselves, define what a particular application should
consider a sufficiently recent valid telemetry sample. The prototype therefore
evaluates its own acceptance and timeout rules.

Failure-detector theory distinguishes completeness and accuracy, including the
possibility that a detector makes mistakes \cite{chandra1996}. That distinction
helps interpret a timeout as a receiver observation of missing accepted data.
It cannot identify which physical component failed, and no formal
failure-detector class is claimed here.

Receiver intervals use Python's \code{time.monotonic()}. Its documented reference
point is unspecified, while differences between readings are meaningful and
unaffected by system-clock updates \cite{python_time}. Sender timestamps remain
recording metadata. Local arrival liveness is therefore separate from a future
model of measurement age or clock synchronization.

\section{System model and liveness rule}
\subsection{Accepted arrivals}
Consider one configured robot and a sequence of accepted receiver arrival times
$r_0,r_1,\ldots$ on the same monotonic clock. The sequence is non-decreasing.
For observation time $t$, define
\begin{equation}
\mathcal{R}(t)=\{r_i:r_i\leq t\}.
\label{eq:arrivals}
\end{equation}
This set contains arrivals that have passed the application checks. In the
default scenario, the last accepted arrival before the blackout is
$r_i=1.9\,\mathrm{s}$. A malformed payload received later does not enter this set.

When $\mathcal{R}(t)$ is nonempty, the most recent accepted arrival has age
\begin{equation}
a(t)=t-\max\mathcal{R}(t).
\label{eq:age}
\end{equation}
The value is an interval in seconds, rather than the age of the underlying
sensor measurement. At $t=2.9\,\mathrm{s}$, with a last arrival at
$1.9\,\mathrm{s}$, Equation~\eqref{eq:age} gives $a(t)=1.0\,\mathrm{s}$.

For timeout $\tau>0$, the observer state is
\begin{equation}
s(t)=
\begin{cases}
\waiting,&\mathcal{R}(t)=\varnothing,\\
\healthy,&\mathcal{R}(t)\neq\varnothing\ \text{and}\ a(t)<\tau,\\
\stale,&\mathcal{R}(t)\neq\varnothing\ \text{and}\ a(t)\geq\tau.
\end{cases}
\label{eq:state}
\end{equation}
With $\tau=1.0\,\mathrm{s}$, the previous example is stale. A new accepted message
at $5.0\,\mathrm{s}$ resets the age to zero and restores the healthy state.
Before the first valid arrival, the observer retains the waiting state.

\subsection{Observation timing}
Suppose no accepted arrival occurs after $r_\ast$, and consecutive observer
evaluations are separated by at most $h$ seconds. The first stale report satisfies
\begin{equation}
r_\ast+\tau\leq t_{\mathrm{detect}}\leq r_\ast+\tau+h.
\label{eq:bound}
\end{equation}
For $r_\ast=1.9\,\mathrm{s}$, $\tau=1.0\,\mathrm{s}$, and an assumed evaluation-gap
bound $h=0.1\,\mathrm{s}$, the interval is $[2.9,3.0]\,\mathrm{s}$. The deterministic
scenario detects the transition at the lower endpoint.

Equation~\eqref{eq:bound} depends on the evaluation-gap assumption. A nominal ROS
timer period does not establish a hard scheduling guarantee on a general operating
system. No real-time bound has been measured. Furthermore, a stream of old
measurements carrying new valid sequence numbers can satisfy arrival liveness
while remaining unsuitable for navigation.

\section{Prototype architecture}
\subsection{Telemetry contract}
The core accepts a versioned JSON object containing robot and publisher-session
identifiers, sequence number, timestamp, geodetic latitude and longitude, heading,
speed, and battery percentage. Units are explicit: degrees, metres per second,
percentage, and timestamp nanoseconds. The source produces illustrative signals,
rather than a physically consistent vessel dynamics model.

RFC~8259 specifies JSON, excludes non-finite numeric values, and discusses
interoperability concerns including duplicate object names \cite{rfc8259}.
The application contract adds strict field and value checks. All expected fields
must be present; extra fields, duplicate keys, and booleans used as numbers are
rejected. Numeric ranges and identifier lengths are bounded.

Payloads are limited to 4096 UTF-8 bytes before JSON decoding. The version-1 speed
limit is $100\,\mathrm{m\,s^{-1}}$, which is an application bound. These checks
bound application parsing and retained state; they do not establish middleware
resource bounds.

\subsection{Sessions and ordering}
Within the active publisher session, accepted sequence numbers must strictly
increase. Gaps are permitted. A previously unseen session permits a publisher
restart with sequence zero. Known retired sessions are rejected, preventing
delayed samples from a known old session from becoming current again.

The monitor retains at most 128 sessions by default. An additional unseen session
is rejected once the limit is reached, while the active session may continue.
Restarting the monitor clears the history. Persistent replay protection and
authenticated session negotiation remain future work.

The model assumes one legitimate publisher for the configured robot identity.
It does not authenticate sources or resolve competing publishers. Invalid
payloads, unexpected robots, duplicate or out-of-order sequences, retired sessions,
and excess new sessions increment a rejection counter without refreshing liveness.

\subsection{Transport separation}
The validator, simulator, monitor, and local demonstration use only Python's
standard library. The default demonstration passes JSON directly between source
and observer within one process.

The optional adapter carries that JSON in \code{std\_msgs/msg/String}.
Telemetry endpoints use KEEP\_LAST, depth~1, BEST\_EFFORT, and VOLATILE.
Status uses depth~1, RELIABLE, and VOLATILE. These settings follow documented
ROS~2 semantics \cite{ros_qos}; compatible transport profiles do not establish
measurement freshness.

In ROS mode, the observer is configured to evaluate and publish status every
$0.1\,\mathrm{s}$ using a steady-clock timer, with separate robot namespaces.
A status consumer still needs to detect disappearance of the observer itself.
Version~0.1.0 contains no vehicle-command interface, generated ROS message
definitions, or MAVLink adapter.

\section{Verification method and results}
\subsection{Deterministic experiment}
The scenario uses duration $D=8.0\,\mathrm{s}$, source rate $f=10\,\mathrm{Hz}$,
and receiver timeout $1.0\,\mathrm{s}$. Sampling includes both endpoints:
\begin{equation}
t_k=\frac{k}{f},\qquad k=0,\ldots,\lfloor Df\rfloor.
\label{eq:samples}
\end{equation}
The configured values give indices $0,\ldots,80$, hence 81 generated samples.

The link model suppresses samples for which $2.0\leq t_k<5.0\,\mathrm{s}$.
These are indices 20 through 49, so
\begin{equation}
N_{\mathrm{generated}}=81,\quad
N_{\mathrm{suppressed}}=30,\quad
N_{\mathrm{accepted}}=81-30=51.
\label{eq:counts}
\end{equation}
Suppression occurs before delivery to the observer, explaining why the base
scenario has zero rejected messages despite missing 30 generated samples.

The observed sequence is waiting at $0.0\,\mathrm{s}$ before the first sample,
healthy at $0.0\,\mathrm{s}$ after acceptance, stale at $2.9\,\mathrm{s}$, and healthy
again at $5.0\,\mathrm{s}$ \cite{artifact}. The final pre-blackout sequence is 19
and the first resumed sequence is 50.

Detection is $2.9-2.0=0.9\,\mathrm{s}$ after the configured blackout starts, or
$2.9-1.9=1.0\,\mathrm{s}$ after the last accepted arrival. These intervals use
different reference events. Recovery at the $5.0\,\mathrm{s}$ virtual tick
demonstrates the state rule rather than a measured network recovery latency.

\subsection{Executed software checks}
The recorded run discovers 19 test methods: 18 core methods pass and one ROS graph
test is skipped because \code{rclpy} is unavailable. The core checks cover timeout
boundaries, recovery, malformed payloads, invalid clocks, sequence ordering,
publisher restarts, retired sessions, history bounds, JSON validation, and
deterministic command-line output \cite{artifact}. This is evidence for the
exercised behaviors rather than a complete correctness proof or coverage result.

An offline wheel build succeeds. Inspection confirms the core modules, ROS
adapter, launch file, package manifest, license, and ament marker. The installed
wheel's JSON demonstration also runs from outside the source tree. Python syntax,
the XML package manifest, and CI YAML parse successfully.

The ROS package build, launch execution, DDS graph test, and ROS CI job were not
executed in the authoring environment. The supplied graph test observes healthy
traffic, cancels the publisher timer to model an outage, and resets it to check
recovery when run in a suitable ROS environment. The test definition alone does
not establish that the transport has passed.

\section{Limitations and next evaluation}
The deterministic scenario uses virtual time and in-process delivery. It does not
exercise discovery delay, scheduling, serialization overhead, queueing, bandwidth
constraints, or a physical wireless link. Arrival age and state labels do not
certify measurement capture time, position accuracy, or vehicle condition.

The next evaluation should execute the Jazzy build and graph test, then record
a real ROS launch with the same blackout. Controlled experiments can vary delay,
jitter, outages, and loss while documenting source rates, session handling, and
timeouts. Detection and recovery delays require explicitly defined ground-truth
events; stale indications during source availability should also be reported.
Repeated runs are needed before estimating performance distributions.

Recorded-data replay, a receive-only MAVLink adapter, and companion-computer
deployment are proposed extensions. Multi-robot field operation, marine dynamics,
new transport performance, and middleware replacement have not been demonstrated.
Research novelty requires a broader literature comparison and additional evidence.

\section{Conclusion}
ROS3-Limassol establishes an OpenPELAOS starting point for explicit telemetry
arrival-liveness monitoring. A receiver monotonic clock, strict message acceptance,
and bounded session history make its behavior inspectable. The deterministic
blackout reproduces stale detection and recovery, and 18 core test methods pass.
A successful ROS transport run and a controlled distributed evaluation are the
next evidence needed. This report records a preliminary implementation baseline.

\section*{Idea attribution}
The initial project idea is credited to \textbf{Christos Paparistodemou}, following
the project initiator's attribution. His public ResearchGate profile is
\url{https://www.researchgate.net/profile/Christos-Paparistodemou}.
The link identifies the contributor; the idea credit is supplied by the project
initiator. Project development and this report are presented under OpenPELAOS
by Charalampos Nadiotis.

\section*{Reproducibility and artifact availability}
The software targets Python~3.10 or newer; the authoring run used Python~3.12.14.
The optional transport target is ROS~2 Jazzy on Ubuntu~24.04. For the current repository layout, first run\\
\code{cd robotics/ros3\_limassol}; then:
\begin{quote}\small
\code{python -m ros3\_limassol.demo}\\
\code{python -m ros3\_limassol.demo --json}\\
\code{python -m unittest discover -s tests -v}
\end{quote}
The accompanying implementation archive includes the Apache-2.0 license,
contributor instructions, tests, and validation record. The public repository is
\url{https://github.com/n0711/ROS3-Limassol}; the robotics package is under
\code{robotics/ros3\_limassol}. The editable XeLaTeX manuscript includes its
bibliography. This is a preliminary, non-peer-reviewed technical report.

\textbf{Version and evidence note (25 October 2026).}
Sections 1--7 retain the original 23 October 2026 version-0.1.0 experiment and
validation record. A subsequent software-development snapshot is version~0.3.1 at
commit \code{272a3b09c2271c56325683f8818e5436a377752f}, adding recording/replay,
a live receive-only recorder, and Humble compatibility work. A previously recorded local rerun of
that commit discovered 45 tests: 42 passed and three ROS integration tests were
skipped because \code{rclpy} was unavailable. These counts must not be confused
with the original 18-test baseline. The repository validation record documents
CI evidence separately. This edition adds a custom manuscript licence, attribution details, and
publication metadata. It preserves the reported experiments and adds no field results.

\clearpage
\section*{Copyright and licence}
\label{sec:licence}
\textbf{Civilian-Use Licence (No Military Use), Version 1.0}\\
\textcopyright\ 2026 Charalampos Nadiotis.

\textbf{1. Licensed material and scope.}
These terms apply to the original copyrightable content of this manuscript,
\emph{ROS3-Limassol: An OpenPELAOS Prototype for Telemetry Liveness Monitoring},
ResearchGate preprint edition, 25 October 2026. Referenced third-party works and
separately licensed software are excluded. Only rights held by the copyright
holder are granted. No patent or trademark rights are granted.

\textbf{2. Permitted civilian reuse.}
Subject to these terms, the copyright holder grants worldwide, royalty-free,
non-exclusive permission to reproduce, distribute, communicate, translate, and
adapt the licensed material for civilian purposes, including commercial civilian
activities.

\textbf{3. Attribution and redistribution.}
Retain the author's name, manuscript title, copyright notice, this licence, and
the credit \emph{Initial project idea: Christos Paparistodemou}. Identify changes
and do not imply endorsement by the author or the credited contributor.
Copies and adaptations containing the licensed material must include these
terms and retain the prohibition on military use. Recipients receive permission
under these same terms; a distributor may not grant military-use permission
for the licensed material.

\textbf{4. Military use prohibited.}
The licensed material and adaptations may not be used for military purposes.
A military purpose is an activity intended to support armed forces, military
operations, or military capabilities. This includes research, design,
development, integration, training, simulation, testing, procurement, deployment,
operation, or maintenance of military weapons, targeting systems, unmanned
systems, command-and-control systems, or military intelligence, surveillance,
and reconnaissance systems. The prohibition applies both to direct military
use and to work by civilian organisations or individuals intended for a military
end use.

\textbf{5. Termination.}
Permission terminates automatically for a user who breaches these terms.
That user must stop exercising the permissions granted by this licence.
Reinstatement requires written permission from the copyright holder.

\textbf{6. Exceptions and warranty.}
Nothing in these terms restricts a use that applicable law permits without
copyright permission. The material is provided \emph{as is}, without warranties,
to the extent permitted by law. These terms regulate copyright reuse of the
manuscript; they do not restrict independent implementations or change rights
already granted for separately licensed software.

\textbf{Software licence boundary.}
The accompanying software remains separately licensed under Apache-2.0.
The manuscript's military-use restriction does not amend that software licence.
This manuscript uses the custom terms above.

\section*{Suggested citation}
Nadiotis, C. (2026).
\emph{ROS3-Limassol: An OpenPELAOS Prototype for Telemetry Liveness Monitoring}.
Preliminary technical report, ResearchGate preprint edition, 25 October 2026.
OpenPELAOS. \url{https://github.com/n0711/ROS3-Limassol}.
Initial project idea: Christos Paparistodemou.

\clearpage
\begin{thebibliography}{9}\small\raggedright
\urlstyle{same}
\setlength{\itemsep}{3pt}
\bibitem{artifact}
C.~Nadiotis. \emph{ROS3-Limassol: OpenPELAOS Development Starter}.
Version~0.1.0, source package and \code{docs/VALIDATION.md}, 23 October 2026.
Software artifact, Apache-2.0. Public repository:
\url{https://github.com/n0711/ROS3-Limassol}.
The original report describes version~0.1.0; subsequent development is identified
separately in the publication-preparation note. No archival DOI assigned.
\bibitem{macenski2022}
S.~Macenski, T.~Foote, B.~Gerkey, C.~Lalancette, and W.~Woodall,
Robot Operating System 2: Design, architecture, and uses in the wild.
\emph{Science Robotics}, vol.~7, no.~66, article eabm6074, 2022.
\href{https://doi.org/10.1126/scirobotics.abm6074}{doi:10.1126/scirobotics.abm6074}.
Author manuscript: \url{https://arxiv.org/abs/2211.07752}.
\bibitem{ros_python}
Open Robotics. Writing a simple publisher and subscriber (Python).
\emph{ROS 2 Documentation: Jazzy}.
\url{https://docs.ros.org/en/jazzy/Tutorials/Beginner-Client-Libraries/Writing-A-Simple-Py-Publisher-And-Subscriber.html}.
Accessed 23 October 2026.
\bibitem{ros_qos}
Open Robotics. Quality of Service settings. \emph{ROS 2 Documentation: Jazzy}.
\url{https://docs.ros.org/en/jazzy/Concepts/Intermediate/About-Quality-of-Service-Settings.html}.
Accessed 23 October 2026.
\bibitem{chandra1996}
T.~D.~Chandra and S.~Toueg. Unreliable failure detectors for reliable distributed
systems. \emph{Journal of the ACM}, vol.~43, no.~2, pp.~225--267, 1996.
\href{https://doi.org/10.1145/226643.226647}{doi:10.1145/226643.226647}.
Author-maintained record: \url{https://www.cs.cornell.edu/info/people/sam/FDpapers.html}.
\bibitem{python_time}
Python Software Foundation. Time access and conversions:
\code{time.monotonic()}. \emph{Python 3.12 Documentation}.
\url{https://docs.python.org/3.12/library/time.html#time.monotonic}.
Accessed 23 October 2026.
\bibitem{rfc8259}
T.~Bray, Ed. \emph{The JavaScript Object Notation (JSON) Data Interchange Format}.
RFC~8259, December 2017.
\url{https://www.rfc-editor.org/rfc/rfc8259.html}.
\end{thebibliography}
\end{document}

